\chapter{Simulacro 4: compresibilidad, ondas de choque y perfiles}
\caja{Condiciones editoriales}{Tiempo sugerido: 110 minutos. Puntuación: 10 puntos. Se permite una calculadora. Esta propuesta no describe las condiciones oficiales de examen. Use ISA con altitud geopotencial, $g_0=9.80665\,\mathrm{m/s^2}$, $R_a=287.05287\,\mathrm{J/(kg\,K)}$ y $\gamma=1.4$.}
\section{Cuestiones razonadas: 2 puntos}
Explique por qué a través de un choque normal estacionario aumenta la entropía y disminuye la presión total mientras la temperatura total se conserva. Distinga el ángulo de Mach del ángulo de un choque oblicuo e indique cuándo coinciden. Justifique por qué la corrección de Prandtl-Glauert no sirve en régimen supersónico ni en las cercanías de $M=1$. Reparta la respuesta en tres explicaciones breves con sus hipótesis. (0.7 + 0.7 + 0.6)
\section{Pitot en vuelo supersónico: 3 puntos}
Una aeronave de investigación vuela en ISA a $H=12000\,\mathrm m$ con $M=2.0$. Se supone gas perfecto, un choque normal desprendido frente a la toma de un tubo Pitot y desaceleración isentrópica posterior hasta la toma.
\begin{enumerate}[label=\alph*)]
\item Calcule $p$, $T$, $\rho$, la velocidad del sonido y la velocidad de vuelo. (0.8)
\item Obtenga el Mach tras el choque y los cocientes de presión, densidad y temperatura estáticas a través de él, además de $p_{t2}/p_{t1}$. (1.2)
\item Calcule la presión de remanso $p_{t2}$ que mediría el Pitot. Si un observador aplicase por error la fórmula subsónica a $p_{t2}$ y $p$, ¿qué Mach leería? Interprete la diferencia. (1.0)
\end{enumerate}
\section{Cuña y choque oblicuo: 3 puntos}
Una corriente de aire con $M_1=2.5$ incide sobre una cuña que produce un choque oblicuo unido de ángulo $\beta=40^\circ$.
\begin{enumerate}[label=\alph*)]
\item Determine el ángulo de desviación $\theta$ de la cuña mediante la relación $\theta$-$\beta$-$M$. (1.0)
\item Calcule los cocientes de presión, densidad y temperatura estáticas a través del choque. (1.0)
\item Obtenga el Mach tras el choque, compruebe que sigue siendo supersónico y explique por qué no existe solución con choque unido si $\theta$ supera un máximo que depende de $M_1$. (1.0)
\end{enumerate}
\section{Perfil y compresibilidad: 2 puntos}
Un perfil de cuerda $c=1.8\,\mathrm m$ vuela a $H=3000\,\mathrm m$ en ISA con $M=0.6$. Use teoría de perfil delgado con pendiente $2\pi\,\mathrm{rad^{-1}}$, $\alpha=4^\circ$ y $\alpha_{L0}=-1.5^\circ$. La viscosidad sigue la ley de Sutherland: $\mu=1.716\times10^{-5}(T/273.15)^{3/2}(273.15+110.4)/(T+110.4)$.
\begin{enumerate}[label=\alph*)]
\item Calcule velocidad de vuelo, densidad y número de Reynolds. (1.0)
\item Obtenga $C_l$ incompresible y con corrección de Prandtl-Glauert; comente su validez. (1.0)
\end{enumerate}

\chapter{Simulacro 5: motores de explosión, de reacción y cohete}
\caja{Condiciones editoriales}{Tiempo sugerido: 120 minutos. Puntuación: 10 puntos. Se permite una calculadora. Esta propuesta no describe las condiciones oficiales de examen. Salvo indicación, use $\gamma=1.4$ y $R_a=287.05287\,\mathrm{J/(kg\,K)}$ con calores específicos constantes. $g_0=9.80665\,\mathrm{m/s^2}$.}
\section{Cuestiones razonadas: 2 puntos}
Explique por qué el rendimiento térmico del ciclo Otto ideal depende solo de la relación de compresión y por qué no puede aumentarse indefinidamente en un motor real. Justifique por qué un turborreactor tiene mayor rendimiento propulsivo cuanto más se acerca la velocidad de salida a la de vuelo y qué se pierde a cambio. Defina impulso específico e indique por qué el valor en vacío es mayor que a nivel del mar para una misma tobera. (0.7 + 0.7 + 0.6)
\section{Ciclo Otto ideal: 2.5 puntos}
Un motor de encendido por chispa admite aire a $T_1=300\,\mathrm K$ y $p_1=95\,\mathrm{kPa}$, con relación de compresión $r=9$. El aporte de calor a volumen constante es $q=1.2\,\mathrm{MJ/kg}$.
\begin{enumerate}[label=\alph*)]
\item Calcule temperatura y presión al final de la compresión isentrópica. (0.8)
\item Calcule temperatura y presión tras el aporte de calor y la temperatura al final de la expansión. (0.9)
\item Determine el rendimiento térmico, el trabajo neto específico y compruebe el balance de calor cedido de dos formas. (0.8)
\end{enumerate}
\section{Turborreactor con tobera adaptada: 2.5 puntos}
Un turborreactor ingiere $45\,\mathrm{kg/s}$ de aire en vuelo a $220\,\mathrm{m/s}$. El dosado másico es $f=0.015$ y los gases salen a $620\,\mathrm{m/s}$ con la presión ambiente.
\begin{enumerate}[label=\alph*)]
\item Calcule el empuje. (1.0)
\item Obtenga la potencia propulsiva, el incremento de energía cinética de la corriente y el rendimiento propulsivo. (1.0)
\item Calcule el consumo específico de empuje en $\mathrm{kg/(N\,s)}$ y en $\mathrm{kg/(N\,h)}$. (0.5)
\end{enumerate}
\section{Tobera de un motor cohete: 3 puntos}
La cámara de un motor cohete está a $p_c=5\,\mathrm{MPa}$ y $T_c=3400\,\mathrm K$. Los gases tienen $\gamma=1.2$ y $R=380\,\mathrm{J/(kg\,K)}$. La garganta tiene $A_t=0.04\,\mathrm{m^2}$ y la tobera se diseña para $p_e=60\,\mathrm{kPa}$. Flujo isentrópico unidimensional.
\begin{enumerate}[label=\alph*)]
\item Calcule el gasto másico, la temperatura y la velocidad de salida. (1.0)
\item Obtenga el área de salida, la relación de expansión y el Mach de salida. (1.0)
\item Calcule el empuje al nivel del mar ($p_a=101325\,\mathrm{Pa}$) y en vacío, y los impulsos específicos correspondientes. Indique el régimen de funcionamiento de la tobera al nivel del mar. (1.0)
\end{enumerate}

\chapter{Simulacro 6: estructuras, materiales y cargas de vuelo}
\caja{Condiciones editoriales}{Tiempo sugerido: 120 minutos. Puntuación: 10 puntos. Se permite una calculadora. Esta propuesta no describe las condiciones oficiales de examen. Los datos de material son docentes: $E=70\,\mathrm{GPa}$ y límite elástico $\sigma_y=300\,\mathrm{MPa}$. Use ISA, $g_0=9.80665\,\mathrm{m/s^2}$ y $R_a=287.05287\,\mathrm{J/(kg\,K)}$.}
\section{Cuestiones razonadas: 2 puntos}
Explique por qué en un recipiente cilíndrico de pared delgada la tensión circunferencial es el doble que la longitudinal. Indique cuándo es válida la fórmula de Euler para una columna y qué ocurre con la carga crítica si se duplica la longitud o se empotran ambos extremos. Defina la velocidad de maniobra y explique qué protege. (0.7 + 0.7 + 0.6)
\section{Fuselaje presurizado: 3 puntos}
Un fuselaje cilíndrico de radio $r=1.9\,\mathrm m$ y espesor $t=1.6\,\mathrm{mm}$ mantiene una altitud de cabina de $2400\,\mathrm m$ mientras la aeronave cruza a $11000\,\mathrm m$ (ISA para ambas).
\begin{enumerate}[label=\alph*)]
\item Calcule las presiones de cabina y exterior y su diferencia. (0.8)
\item Obtenga las tensiones circunferencial y longitudinal, la tensión equivalente de von Mises y el coeficiente de seguridad estático. (1.2)
\item Explique por qué ese coeficiente no garantiza la vida a fatiga ni la tolerancia al daño. (1.0)
\end{enumerate}
\section{Pandeo de una barra tubular: 2.5 puntos}
Una barra tubular biarticulada de diámetro exterior $40\,\mathrm{mm}$, interior $36\,\mathrm{mm}$ y longitud $1.2\,\mathrm m$ trabaja a compresión.
\begin{enumerate}[label=\alph*)]
\item Calcule momento de inercia, área y carga crítica de Euler. (1.0)
\item Obtenga la esbeltez, compruebe la validez de Euler con la esbeltez límite $\pi\sqrt{E/\sigma_y}$ y determine si gobierna el pandeo o la fluencia. (1.0)
\item Indique cómo cambia la carga crítica con empotramiento doble. (0.5)
\end{enumerate}
\section{Maniobra y ráfaga: 2.5 puntos}
Una aeronave ligera de $W=11800\,\mathrm N$ y $S=16\,\mathrm{m^2}$ tiene $C_{L\max}=1.5$. Se toma $\rho_0=1.225\,\mathrm{kg/m^3}$ y un factor de carga límite $n_{\max}=3.8$.
\begin{enumerate}[label=\alph*)]
\item Calcule la velocidad de pérdida con $n=1$ y la velocidad de maniobra. (1.0)
\item Para una ráfaga discreta de $7.6\,\mathrm{m/s}$ a una velocidad equivalente de $45\,\mathrm{m/s}$ y pendiente de sustentación $5.0\,\mathrm{rad^{-1}}$, calcule el incremento y el factor de carga resultante, sin factor de atenuación. (1.0)
\item Compare con el límite e indique dos simplificaciones del modelo. (0.5)
\end{enumerate}

\chapter{Simulacro 7: navegación, ortodrómica y actuaciones de pista}
\caja{Condiciones editoriales}{Tiempo sugerido: 110 minutos. Puntuación: 10 puntos. Se permite una calculadora. Esta propuesta no describe las condiciones oficiales de examen. Tierra esférica de radio $6371\,\mathrm{km}$; ISA con altitud geopotencial, $g_0=9.80665\,\mathrm{m/s^2}$, $R_a=287.05287\,\mathrm{J/(kg\,K)}$.}
\section{Cuestiones razonadas: 2 puntos}
Distinga velocidad verdadera, velocidad respecto al suelo y velocidad indicada. Explique por qué una ruta ortodrómica no mantiene un rumbo constante y cuándo conviene la loxodrómica. Justifique por qué un aeródromo elevado en un día cálido alarga el recorrido de despegue. (0.7 + 0.7 + 0.6)
\section{Triángulo de vientos: 3 puntos}
Una aeronave de $\mathrm{TAS}=75\,\mathrm{m/s}$ debe seguir una derrota $090^\circ$. El viento sopla de $030^\circ$ con $15\,\mathrm{m/s}$.
\begin{enumerate}[label=\alph*)]
\item Calcule las componentes longitudinal y transversal del viento respecto a la derrota. (0.8)
\item Determine el ángulo de corrección de deriva, el rumbo a volar y la velocidad respecto al suelo. (1.4)
\item Calcule el tiempo de un tramo de $180\,\mathrm{km}$. (0.8)
\end{enumerate}
\section{Distancia y rumbo inicial ortodrómicos: 2.5 puntos}
Entre dos aeropuertos de coordenadas redondeadas $40.47^\circ\mathrm N,\,3.56^\circ\mathrm W$ y $41.30^\circ\mathrm N,\,2.08^\circ\mathrm E$.
\begin{enumerate}[label=\alph*)]
\item Calcule el ángulo central por la fórmula haversine y la distancia ortodrómica. (1.5)
\item Obtenga el rumbo inicial y explique su variación a lo largo de la ruta. (1.0)
\end{enumerate}
\section{Recorrido de despegue: 2.5 puntos}
Una aeronave ligera de $W=15000\,\mathrm N$, $S=16\,\mathrm{m^2}$ y $C_{L\max}=1.6$ despega con empuje estático constante $3200\,\mathrm N$ al nivel del mar ($\rho_0=1.225\,\mathrm{kg/m^3}$), con coeficiente de rodadura $\mu=0.03$. Se modela una aceleración media constante $a=g_0(T/W-\mu)$, sin resistencia aerodinámica, y la velocidad de despegue es $1.2V_s$.
\begin{enumerate}[label=\alph*)]
\item Calcule el recorrido al nivel del mar. (0.8)
\item Repítalo en un aeródromo a $2000\,\mathrm m$ ISA con temperatura ISA$+20\,\mathrm K$, suponiendo que el empuje es proporcional a la densidad. (1.2)
\item Compare ambos resultados e indique tres efectos que el modelo omite. (0.5)
\end{enumerate}

\chapter{Simulacro 8: sistemas hidráulicos, eléctricos y alcance}
\caja{Condiciones editoriales}{Tiempo sugerido: 110 minutos. Puntuación: 10 puntos. Se permite una calculadora. Esta propuesta no describe las condiciones oficiales de examen. Datos educativos; no representan una aeronave ni un componente certificados. $g_0=9.80665\,\mathrm{m/s^2}$.}
\section{Cuestiones razonadas: 2 puntos}
Explique por qué un actuador hidráulico de doble efecto con vástago de un solo lado desarrolla menos fuerza y más velocidad en retracción que en extensión. Distinga la capacidad nominal de una batería de la energía realmente utilizable. Justifique por qué la ecuación de Breguet supone cociente $L/D$ y consumo específico constantes y qué ocurre si no lo son. (0.7 + 0.7 + 0.6)
\section{Actuador hidráulico: 3 puntos}
Un actuador de doble efecto tiene pistón de $50\,\mathrm{mm}$ de diámetro y vástago de $25\,\mathrm{mm}$. El sistema trabaja a $21\,\mathrm{MPa}$ y el caudal disponible es $1.0\times10^{-3}\,\mathrm{m^3/s}$.
\begin{enumerate}[label=\alph*)]
\item Calcule la fuerza máxima en extensión y en retracción. (0.8)
\item Calcule las velocidades de extensión y de retracción con ese caudal. (0.7)
\item Con una carga de $30\,\mathrm{kN}$ en extensión, calcule potencia hidráulica, potencia útil y potencia disipada; la potencia de entrada de una bomba de rendimiento $0.85$; y la tasa de calentamiento si $18\,\mathrm{kg}$ de fluido ($c=1900\,\mathrm{J/(kg\,K)}$) absorben el calor sin enfriamiento. (1.5)
\end{enumerate}
\section{Batería y cableado: 2.5 puntos}
Una batería de $24\,\mathrm V$ y $40\,\mathrm{Ah}$ alimenta una carga de $15\,\mathrm A$ continuos mediante un cable de cobre ($\rho_{Cu}=1.72\times10^{-8}\,\Omega\,\mathrm m$) de $6\,\mathrm{mm^2}$ y $12\,\mathrm m$ de longitud total de ida y vuelta.
\begin{enumerate}[label=\alph*)]
\item Calcule la autonomía si solo es utilizable el $80\,\%$ de la capacidad. (0.8)
\item Calcule la resistencia del cable, la caída de tensión y la pérdida por efecto Joule. (1.2)
\item Indique la tensión en la carga y la caída relativa. (0.5)
\end{enumerate}
\section{Alcance y autonomía de Breguet: 2.5 puntos}
Una aeronave a reacción vuela a $230\,\mathrm{m/s}$ con $L/D=16$ y consumo específico $1.7\times10^{-5}\,\mathrm{kg/(N\,s)}$. La razón entre pesos inicial y final es $1.25$.
\begin{enumerate}[label=\alph*)]
\item Calcule el alcance de Breguet. (1.2)
\item Calcule la autonomía. (0.8)
\item Obtenga la fracción de combustible respecto al peso inicial e indique dos condiciones que la ecuación supone. (0.5)
\end{enumerate}
